% 正文主体（由 main.tex / main_blind.tex 在设置 \ifblindthesis 后 \input）
% geometry 必须早于 fancyhdr 加载，否则页眉页脚居中会相对于版心而非纸张
\usepackage{geometry}
% a4paper：与国内 Word 默认及硕论常见要求一致，避免 Letter(612x792) 与封面 PDF 拼版时版心不一致
% 上边距：减小 top + headsep，避免“页眉实际用 eso-pic 画、版心仍留一大块 head 区域”导致正文起点过低
\geometry{a4paper, left=2.5cm, right=2.5cm, top=22mm, bottom=2.5cm, headheight=10pt, headsep=8mm}
\input{images/option}
\renewcommand{\textfraction}{0.15}    % 正文至少占页面的15%（避免图片占满整页）
% \renewcommand{\floatpagefraction}{0.8}% 图片占页面80%以上才会单独占页
\renewcommand{\topfraction}{0.85}     % 顶部浮动体最多占页面85%
\renewcommand{\bottomfraction}{0.85}

\addbibresource{chapters/reference.bib}

% --- 1. 基础宏包 ---
\usepackage{amsmath, amssymb, amsfonts} % 数学公式 (去重)
\usepackage{indentfirst}
\usepackage{graphicx}   % 图片
\usepackage{booktabs}   % 三线表
\usepackage{array}      % 表格列格式（如 >{...}p{...}）
\usepackage{multirow}   % 表格跨行
\usepackage{fancyhdr}   % 页眉页脚（geometry 已在文件开头加载）
\usepackage{eso-pic}    % 用于在纸张坐标系中绝对居中绘制页眉
\usepackage{hyperref}   % 超链接
\usepackage{titletoc}   % 目录格式
\usepackage{enumitem}
% 中文正文列表：一级为全角「（1）（2）…」+ 悬挂缩进；二级为「1. 2. …」（用于条目下再分点）
\setlist[enumerate,1]{
  label=（\arabic*）,
  leftmargin=2.85em,
  labelwidth=2.25em,
  labelsep=0.35em,
  align=left,
  itemsep=0.12\baselineskip,
  parsep=0pt,
  topsep=0.22\baselineskip,
  partopsep=0pt,
}
\setlist[enumerate,2]{
  label=\arabic*.,
  leftmargin=3.45em,
  labelsep=0.4em,
  align=left,
  itemsep=0.06\baselineskip,
  parsep=0pt,
  topsep=0.1\baselineskip,
}
\setlist[itemize,1]{
  leftmargin=2.4em,
  itemsep=0.1\baselineskip,
  parsep=0pt,
  topsep=0.2\baselineskip,
}
% 中文「（1）（2）…」：首行（含编号）整体 2em 缩进，续行顶格；编号「）」后 1 空格（非 enumerate 悬挂）
\newcounter{cnparenum}
\newenvironment{cnenumerate}{%
  \begingroup
  \setlength{\parindent}{0pt}%
  \setcounter{cnparenum}{0}%
  \vspace{0.22\baselineskip}%
}{%
  \par\vspace{0.22\baselineskip}%
  \endgroup
}
\newcommand{\cnitem}{%
  \ifnum\value{cnparenum}>0
    \par\addvspace{0.12\baselineskip}%
  \else
    \par
  \fi
  \refstepcounter{cnparenum}%
  \parshape 2 2em \dimexpr\linewidth-2em\relax 0pt \linewidth
  （\arabic{cnparenum}）\space
  \ignorespaces
}
% 套在 cnenumerate 内的二级列表（因外层非 list，需显式指定 1. 2. …）
\newenvironment{subenumerate}{%
  \begin{enumerate}[
    label=\arabic*.,
    leftmargin=3.45em,
    labelsep=0.4em,
    align=left,
    itemsep=0.06\baselineskip,
    parsep=0pt,
    topsep=0.1\baselineskip,
    partopsep=0pt,
  ]}%
  {\end{enumerate}}
\usepackage{algorithm}    % 算法浮动环境
\usepackage{algorithmicx} % 算法结构扩展
\usepackage{algpseudocode}% 算法关键词（For/If/State等）
% 算法题注：与图/表一致（图/表/算法 与编号之间无空格）；名称中文
\floatname{algorithm}{算法}
\captionsetup[algorithm]{labelformat=nwunosep,labelsep=quad,labelfont={nwulabel},font={nwucaption},skip=5pt,belowskip=3pt}
\usepackage{amsmath}      % 数学公式支持
\usepackage{amssymb}
% 注意：参考文献使用 biblatex（在 option.tex 中已配置），不需要 natbib

% --- 2. 算法宏包 (修复冲突) ---
% \usepackage{algorithmic} % 已注释：与 algpseudocode 冲突，建议不用


% --- 3. 页眉页脚（西北大学：从绪论开始）---
% 分析：fancyhdr 的 [C] 以 typeblock 为基准，\headwidth=\paperwidth 无法改变起始位置，故无法相对纸张居中。
% 方案：用 eso-pic 在页顶附近绘制页眉；竖向位置与 \geometry{top, headheight, headsep} 大致匹配，避免版心顶端空白过高
\newcommand{\headerfont}{\xiaowuhao\songti}
\newcounter{usefancyheader}
\pagestyle{empty}
% 前置部分（摘要、目录）页码：小五号 Times New Roman 罗马数字
\fancypagestyle{front}{%
  \fancyhf{}
  \renewcommand{\headrulewidth}{0pt}
  \fancyfoot[C]{\wuhao\rmfamily\thepage}
}
% eso-pic：仅在主体部分（绪论起）绘制页眉，与版心同宽、同左边界，使页眉与正文标题水平对齐
% 策略：版心左边界为 1in+\oddsidemargin（与 geometry left=3cm 一致），宽度 \textwidth；
%       页眉文字与页眉线均置于该区域内并居中，与各章标题、致谢等保持同一中轴
\AddToShipoutPictureFG{%
  \ifnum\value{usefancyheader}=1
    \AtPageUpperLeft{%
      \put(\dimexpr1in+\oddsidemargin\relax,-\dimexpr9mm\relax){\makebox[\textwidth][c]{\headerfont\ifodd\value{page}\leftmark\else 西北大学硕士学位论文\fi}}%
      \put(\dimexpr1in+\oddsidemargin\relax,-\dimexpr14mm\relax){\rule{\textwidth}{0.4pt}}%
    }%
  \fi
}
% fancy/plain：页眉文字和线由 eso-pic 绘制，此处仅保留页脚
\fancypagestyle{fancy}{%
  \fancyhf{}
  % oneside 下不区分奇偶页，避免使用 E/O 选项触发 fancyhdr 警告
  \fancyhead{}
  \renewcommand{\headrulewidth}{0pt}
  % oneside 下不区分奇偶页，用 C 保持居中即可
  \fancyfoot[C]{\wuhao\rmfamily\thepage}
}
\fancypagestyle{plain}{%
  \fancyhf{}
  % oneside 下不区分奇偶页，避免使用 E/O 选项触发 fancyhdr 警告
  \fancyhead{}
  \renewcommand{\headrulewidth}{0pt}
  \fancyfoot[C]{\wuhao\rmfamily\thepage}
}
% 章间仅换页，不插入空白页（单面印刷可减少大段留白；若学院要求章从奇数页起，可恢复奇数页逻辑）
\newcommand{\chapterblank}{\clearpage}
% ctexart 无 \mainmatter，定义为空操作以免报错
\providecommand{\mainmatter}{}

% --- 信息录入 ---
\title{基于对比学习的联邦异构图推荐算法与系统研究}
\author{伍勋高}
\date{2026年3月}

\begin{document}
\raggedbottom

% =========================================================
%   前置部分 (Front Matter)
%   特点：页码为罗马数字 (i, ii, iii...)
% =========================================================

% 从摘要开始启用页眉（eso-pic 绘制）
\setcounter{usefancyheader}{1}

% 1. 摘要页 (中英文摘要)
\input{chapters/abstract}

% 4. 目录（西北大学规范：一般为二级或三级）
\setcounter{tocdepth}{3}
\newpage
\markboth{目\quad 录}{目\quad 录}
\par\addvspace{1.5em}
\begin{center}
\xiaoerhao\heiti 目~~~~录
\end{center}
\par\addvspace{1em}
\begingroup
  % 目录不用正文 1.5 倍行距，否则条目标题行间显稀
  \renewcommand{\baselinestretch}{1}\selectfont
  \StartToc{toc}
\endgroup
\newpage  % 目录页（包含 \tableofcontents）
\newpage
% 图目录（在总目录中以“图索引”形式出现，层级与章节标题一致）
\markboth{图\quad 目\quad 录}{图\quad 目\quad 录}
\phantomsection
\addcontentsline{toc}{section}{图索引}
\begin{center}
\sanhao\heiti 图~~~~目~~~~录
\end{center}
\addvspace{1em}
\StartToc{lof}
\newpage
% 表目录（在总目录中以“表索引”形式出现，层级与章节标题一致）
\markboth{表\quad 目\quad 录}{表\quad 目\quad 录}
\phantomsection
\addcontentsline{toc}{section}{表索引}
\begin{center}
\sanhao\heiti 表~~~~目~~~~录
\end{center}
\addvspace{1em}
\StartToc{lot}
\newpage

% =========================================================
%   主体部分 (Main Matter)
%   特点：页码重置为阿拉伯数字 (1, 2, 3...)；绪论从奇数页起，页眉从绪论开始
% =========================================================
\mainmatter
\clearpage
% 主体部分保持启用页眉
\input{chapters/chap1_intro}
\chapterblank
\input{chapters/chap2_related}
\chapterblank
\input{chapters/chap3_method}
\chapterblank
\input{chapters/chap4_compression}
\chapterblank
\input{chapters/chap5_system}
\chapterblank
\input{chapters/chap6_summary}
\chapterblank

% =========================================================
%   后置部分 (Back Matter)
%   特点：参考文献、致谢、成果等
% =========================================================
% \backmatter  % ctexart 不支持 \backmatter，已注释

% 1. 参考文献
% 使用 biblatex（已在 option.tex 中配置）
% Reference.tex 中已包含 \printbibliography
\newpage
\markboth{参考文献}{参考文献}
% 参考文献在目录中与“摘要”“第一章 绪论”等保持同一层级与样式（section 级）
\addcontentsline{toc}{section}{参考文献}
% 参考文献标题格式：三号黑体，居中
\begin{center}
{\fontsize{15.75pt}{\baselineskip}\selectfont\heiti\bfseries 参考文献}
\end{center}
\vspace{0.5cm}
% 参考文献内容：五号宋体
\printbibliography[heading=none]  % 不显示默认标题，使用自定义标题
\newpage

% 2. 致谢 (保留一个即可)
\input{chapters/thank}
% \input{backmatter/ack} % 已注释：避免重复致谢

% 3. 攻读硕士学位期间取得的科研成果
% 按要求删除该部分
% % 攻读硕士学位期间取得的科研成果（与参考文献、致谢同级目录）
\newpage
\markboth{攻读硕士学位期间取得的科研成果}{攻读硕士学位期间取得的科研成果}
\phantomsection
\addcontentsline{toc}{section}{攻读硕士学位期间取得的科研成果}

\begin{center}
{\fontsize{15.75pt}{\baselineskip}\selectfont\heiti 攻读硕士学位期间取得的科研成果}
\end{center}
\vspace{0.35\baselineskip}

% 与中文摘要、正文一致：小四号宋体，行距 12pt/18pt（约 1.5 倍），西文与数字随正文主字体（Times New Roman）
\setlength{\parindent}{2em}
\setlength{\parskip}{0pt}
\begingroup
\fontsize{12pt}{18pt}\selectfont

% 不用 \subsubsection*，避免 titlesec 块级标题带来大块竖向留白
\noindent\heiti 参与科研项目及科研获奖\par
\vspace{0.35\baselineskip}
\setlength{\parindent}{0pt}
\songti
\begin{enumerate}[
  label={[\arabic*]},
  align=left,
  leftmargin=*,
  itemindent=0pt,
  labelindent=0pt,
  labelwidth=1.35em,
  labelsep=0.15em,
  itemsep=0.15\baselineskip,
  parsep=0pt,
  topsep=0pt,
  partopsep=0pt,
]
  \item 校级项目，西北大学 IOC 智慧运营中心建设项目，2024.09--2025.06，项目组主要成员。
  \item 校级项目，西北大学导师上岗系统，2024.10--2024.12，项目组主要成员。
\end{enumerate}
\endgroup

% 文末独立空白页（双面印刷时常用；单面下亦保留版式空白）
\newpage
\thispagestyle{empty}
\null
\newpage


% \input{} % 已删除：空的 input 会导致编译错误

\end{document}
